%% ma: L10-B14-0074
%% lop: 10
%% chuong: 5
%% bai: 14
%% dang: 10.14.2
%% loai: TN4
%% muc: TH
%% dap_an: "A"
%% nguon: "chuyen-de-cac-so-dac-trung-cua-mau-so-lieu-khong-ghep-nhom-toan-10-chuong-trinh-moi.pdf (GV. Phan Nhật Linh) (P05), Bài 03 Dạng 2 Phần I Câu 21"
%% kien_thuc: "Lớp 10 Bài 13 (tứ phân vị), Bài 14 (khoảng tứ phân vị)"
%% trang_thai: chinh-thuc
%% ghi_chu: "Cùng số liệu Bài 03 Dạng 1 Bài tập tự luận 8 nhưng khác loại câu. Tự duyệt theo yêu cầu giáo viên 22/10/2026."
\begin{ex}
Số lượng vải bán ra tại một quầy hàng các tháng trong một năm được cho bởi bảng sau (đơn vị: mét):
\begin{center}
\begin{tabular}{|c|c|c|c|c|c|c|c|c|c|c|c|c|}\hline
Tháng & $1$ & $2$ & $3$ & $4$ & $5$ & $6$ & $7$ & $8$ & $9$ & $10$ & $11$ & $12$\\\hline
Số lượng & $410$ & $430$ & $450$ & $430$ & $525$ & $550$ & $950$ & $450$ & $800$ & $635$ & $760$ & $560$\\\hline
\end{tabular}
\end{center}
Khoảng tứ phân vị của mẫu số liệu là
\choice{\True $\Delta_Q=257{,}5$}{$\Delta_Q=255{,}5$}{$\Delta_Q=257$}{$\Delta_Q=255$}
\loigiai{Sắp xếp: $410$; $430$; $430$; $450$; $450$; $525$; $550$; $560$; $635$; $760$; $800$; $950$.

$Q_1=440$; $Q_3=697{,}5$; $\Delta_Q=257{,}5$. Chọn A.}
\end{ex}
